\documentclass[UTF8, 12pt, fontset=fandol, oneside]{ctexart}
\usepackage{researchnotes}

\graphicspath{{math/概率论/}}
\newcommand{\biohazard}{\ensuremath{\ast}}

\title{第 1 章\quad 概率与计数}
\author{}
\date{}

\begin{document}

\maketitle

\noindent\textit{运气、巧合、随机性、不确定性、风险、怀疑、命运、机会。}
这些词我们大概已经听过无数次，但它们常常只是以含混而随意的方式出现。概率在科学和日常生活中无处不在，但它也可能极其违反直觉。如果依赖可靠性可疑的直觉，我们就很容易作出不准确的预测，或在决策中过度自信。本书的目标，是把概率作为一种逻辑框架来介绍，用有原则的方式量化不确定性和随机性。同时，我们也会努力增强直觉：既包括直觉与逻辑推理相一致的情形，也包括直觉并不那么幸运的情形。

\section*{1.1\quad 为什么学习概率？}

数学是关于确定性的逻辑，而概率是关于不确定性的逻辑。概率在许多领域都极其有用，因为它提供了理解和解释变异、从噪声中分离信号、以及刻画复杂现象的工具。下面只是从不断增长的应用清单中取出的一个很小样本：

\begin{enumerate}
  \item \textit{统计学：}概率是统计学的基础和语言，使得许多强有力的方法能够利用数据来认识世界。
  \item \textit{物理学：}爱因斯坦有一句著名的话：“上帝不掷骰子。”但人们今天对量子物理的理解，在自然界最基本的层面上大量涉及概率。统计力学也是建立在概率之上的物理学重要分支。
  \item \textit{生物学：}遗传学与概率紧密交织，无论是在基因遗传中，还是在随机突变建模中都是如此。
  \item \textit{计算机科学：}随机算法在运行时作出随机选择；在许多重要应用中，它们比目前已知的确定性替代方法更简单、更高效。概率在研究算法性能、机器学习和人工智能中也起着核心作用。
  \item \textit{气象学：}天气预报是（或者应该是）用概率来计算和表达的。
  \item \textit{赌博：}概率的许多早期研究，都是为了回答与赌博和机会游戏有关的问题。
  \item \textit{金融学：}即使与前一个例子略有重复，也仍然值得指出，概率在量化金融中居于中心地位。股票价格随时间的建模，以及金融工具“公平”价格的确定，都严重依赖概率。
  \item \textit{政治学：}近年来，政治学越来越定量化、统计化，其应用包括分析民意调查、评估选区划分问题、预测选举等。
  \item \textit{医学：}随机临床试验的发展改变了近年的医学研究。在这类试验中，患者被随机分配接受治疗或安慰剂。生物统计学家 David Harrington 曾评论说，有人猜测这可能是二十世纪科学医学中最重要的进展；现代科学中一个有趣的反讽是，随机试验通过在研究设计中引入机会变异，同时调整了受控实验中已观察和未观察到的异质性。
  \item \textit{生活：}生活充满不确定性，而概率是不确定性的逻辑。虽然不可能对生活中的每一个决定都进行正式的概率计算，但认真思考概率可以帮助我们避免一些常见谬误，解释巧合，并作出更好的预测。
\end{enumerate}

概率提供了有原则地解决问题的程序，但它也可能带来陷阱和悖论。例如，本章中我们会看到，即使是 17 世纪独立发现微积分的两个人 Gottfried Wilhelm von Leibniz 和 Sir Isaac Newton，也不能免于概率中的基本错误。贯穿本书，我们会用以下策略来帮助避免潜在的陷阱。

\begin{enumerate}
  \item \textit{模拟：}概率中一个美妙的方面是，我们常常可以通过模拟来研究问题。与其无休止地和意见不同的人争论答案，不如运行模拟，经验性地看看谁是对的。
  \item \textit{生物危害标记：}研究常见错误很重要，它能帮助我们更深入地理解什么是有效的概率推理、什么不是。在本书中，常见错误称为“生物危害”，并用 $\biohazard$ 标记，因为犯这类错误可能“危害健康”。
  \item \textit{健全性检查：}用一种方法解决问题之后，我们常常会尝试用另一种方法重新求解，或者检查答案在简单情形和极端情形下是否合理。
\end{enumerate}

\section*{1.2\quad 样本空间与鹅卵石世界}

概率的数学框架围绕集合建立。设想进行一次实验，它的结果会是某个可能结果集合中的一个。在实验进行之前，我们不知道哪一个结果会出现；实验之后，结果便“结晶”为实际结果。

\noindent\textbf{定义 1.2.1（样本空间与事件）}\quad 一个实验的\textbf{样本空间} $S$，是该实验所有可能结果组成的集合。一个\textbf{事件} $A$ 是样本空间 $S$ 的一个子集；若实际结果属于 $A$，则称事件 $A$ 发生。

\begin{figure}[htbp]
  \centering
  \includegraphics[width=0.58\textwidth]{1-1.png}
  \caption{作为鹅卵石世界的样本空间，其中突出显示了两个事件 $A$ 和 $B$。}
\end{figure}

一个实验的样本空间可以是有限的、可数无穷的，或者不可数无穷的。当样本空间有限时，我们可以把它想象成图 1.1 所示的\textbf{鹅卵石世界}。每一颗鹅卵石代表一个结果，而一个事件就是一组鹅卵石。

进行实验相当于随机选择一颗鹅卵石。如果所有鹅卵石具有相同的质量，那么每颗鹅卵石被选中的可能性相同。这个特殊情形是接下来两节的主题。在 1.6 节中，我们会给出概率的一般定义，允许鹅卵石具有不同的质量。

集合论在概率中非常有用，因为它为表达和处理事件提供了丰富语言。集合运算，尤其是并、交和补，使我们能够用已经定义好的事件构造新事件。这些概念也使我们能够用不止一种方式表达同一个事件；通常，同一个事件的某种表达方式会比另一种更容易处理。

例如，设 $S$ 是某个实验的样本空间，$A,B\subseteq S$ 是事件。则并集 $A\cup B$ 是当且仅当 $A,B$ 中至少一个发生时发生的事件；交集 $A\cap B$ 是当且仅当 $A$ 与 $B$ 都发生时发生的事件；补集 $A^c$ 是当且仅当 $A$ 不发生时发生的事件。我们还有 De Morgan 律：
\[
  (A\cup B)^c=A^c\cap B^c,\qquad (A\cap B)^c=A^c\cup B^c,
\]
因为说“并非 $A$ 与 $B$ 至少一个发生”，等价于说 $A$ 不发生且 $B$ 不发生；说“并非二者都发生”，等价于说至少有一个不发生。类似结果也适用于两个以上事件的并和交。

在图 1.1 所示的例子中，$A$ 是由 5 颗鹅卵石组成的集合，$B$ 是由 4 颗鹅卵石组成的集合；$A\cup B$ 由属于 $A$ 或 $B$ 的 8 颗鹅卵石组成（包括同时属于二者的那颗），$A\cap B$ 由同时属于 $A$ 与 $B$ 的那颗鹅卵石组成，而 $A^c$ 由不属于 $A$ 的 4 颗鹅卵石组成。

样本空间的概念非常一般而抽象，因此心中有一些具体例子很重要。

\noindent\textbf{例 1.2.2（掷硬币）}\quad 将一枚硬币抛掷 10 次。记正面为 $H$，反面为 $T$，一个可能结果（鹅卵石）是 $HHHTHHTTHT$；样本空间是所有长度为 10、由 $H$ 和 $T$ 组成的字符串的集合。我们可以（也将会）把 $H$ 编码为 $1$，把 $T$ 编码为 $0$，于是一个结果就是序列 $(s_1,\ldots,s_{10})$，其中 $s_j\in\{0,1\}$，样本空间则是所有这类序列的集合。现在看一些事件：

\begin{enumerate}
  \item 令 $A_1$ 为第一次抛掷为正面的事件。作为集合，
  \[
    A_1=\{(1,s_2,\ldots,s_{10}):s_j\in\{0,1\},\ 2\le j\le 10\}.
  \]
  这是样本空间的一个子集，因此它确实是一个事件；说 $A_1$ 发生，等价于说第一次抛掷为正面。类似地，对 $j=2,3,\ldots,10$，令 $A_j$ 为第 $j$ 次抛掷为正面的事件。

  \item 令 $B$ 为至少有一次抛掷为正面的事件。作为集合，
  \[
    B=\bigcup_{j=1}^{10} A_j.
  \]

  \item 令 $C$ 为所有抛掷都为正面的事件。作为集合，
  \[
    C=\bigcap_{j=1}^{10} A_j.
  \]

  \item 令 $D$ 为至少出现两个连续正面的事件。作为集合，
  \[
    D=\bigcup_{j=1}^{9}(A_j\cap A_{j+1}).
  \]
\end{enumerate}

\noindent\textbf{例 1.2.3（任取一张牌）}\quad 从一副标准的 52 张牌中抽取一张牌。样本空间 $S$ 是全部 52 张牌组成的集合（因此有 52 颗鹅卵石，每张牌对应一颗）。考虑下面四个事件：
\begin{itemize}
  \item $A$：抽到 A；
  \item $B$：抽到黑色花色；
  \item $D$：抽到方块；
  \item $H$：抽到红心。
\end{itemize}
作为集合，$H$ 由 13 张牌组成：
\[
  \{\text{红心 A},\text{红心 2},\ldots,\text{红心 K}\}.
\]
我们可以用事件 $A,B,D,H$ 构造许多其他事件。并、交和补对此尤其有用。例如：
\begin{itemize}
  \item $A\cap H$ 是抽到红心 A 的事件；
  \item $A\cap B$ 是事件 $\{\text{黑桃 A},\text{梅花 A}\}$；
  \item $A\cup D\cup H$ 是抽到红牌或 A 的事件；
  \item $(A\cup B)^c=A^c\cap B^c$ 是抽到红色非 A 牌的事件。
\end{itemize}
另外，注意到 $(D\cup H)^c=D^c\cap H^c=B$，所以 $B$ 可以用 $D$ 和 $H$ 来表示。另一方面，抽到黑桃的事件不能只用 $A,B,D,H$ 来写出，因为这几个事件都不够精细，无法区分黑桃和梅花。

在这个样本空间中还可以定义许多其他事件。事实上，本章后面介绍的计数方法会说明，即使只有 52 颗鹅卵石，这个问题中也有
\[
  2^{52}\approx 4.5\times 10^{15}
\]
个事件。

如果抽到的牌是一张 joker 会怎样？这说明我们选错了样本空间；我们假定实验结果一定是 $S$ 的一个元素。

前面的例子说明，事件既可以用自然语言描述，也可以用集合记号描述。有时自然语言描述更容易理解，而集合记号更容易操作。设 $S$ 为样本空间，$s_{\text{actual}}$ 为实验的实际结果（即实验进行时最终被选中的那颗鹅卵石）。下一页给出一个在自然语言和集合之间转换的小词典。

\begin{center}
\begin{tabular}{ll}
\hline
\textbf{自然语言} & \textbf{集合记号} \\
\hline
\multicolumn{2}{l}{\textit{事件与发生}}\\
样本空间 & $S$\\
$s$ 是一个可能结果 & $s\in S$\\
$A$ 是一个事件 & $A\subseteq S$\\
$A$ 发生 & $s_{\text{actual}}\in A$\\
一定会发生某事 & $s_{\text{actual}}\in S$\\
\hline
\multicolumn{2}{l}{\textit{由旧事件构造新事件}}\\
$A$ 或 $B$（包含式） & $A\cup B$\\
$A$ 且 $B$ & $A\cap B$\\
非 $A$ & $A^c$\\
$A$ 或 $B$，但不同时发生 & $(A\cap B^c)\cup(A^c\cap B)$\\
$A_1,\ldots,A_n$ 中至少一个 & $A_1\cup\cdots\cup A_n$\\
$A_1,\ldots,A_n$ 全部 & $A_1\cap\cdots\cap A_n$\\
\hline
\multicolumn{2}{l}{\textit{事件之间的关系}}\\
$A$ 蕴含 $B$ & $A\subseteq B$\\
$A$ 与 $B$ 互斥 & $A\cap B=\varnothing$\\
$A_1,\ldots,A_n$ 是 $S$ 的一个划分 &
$A_1\cup\cdots\cup A_n=S,\ A_i\cap A_j=\varnothing\ (i\ne j)$\\
\hline
\end{tabular}
\end{center}

\section*{1.3\quad 概率的朴素定义}

从历史上看，事件概率最早的定义，是数出该事件可能发生的方式数，再除以实验所有可能结果的总数。我们称它为\textbf{朴素定义}，因为它限制很强，并依赖强假设；不过，只要不误用，它仍然重要而有用。

\noindent\textbf{定义 1.3.1（概率的朴素定义）}\quad 设 $A$ 是一个实验的事件，该实验具有有限样本空间 $S$。事件 $A$ 的\textbf{朴素概率}为
\[
  P_{\text{naive}}(A)=\frac{|A|}{|S|}
  =\frac{\text{有利于 }A\text{ 的结果数}}{S\text{ 中结果总数}}.
\]
这里用 $|A|$ 表示 $A$ 的大小。

用鹅卵石世界来讲，朴素定义只是说，$A$ 的概率就是位于 $A$ 中的鹅卵石所占的比例。例如，在图 1.1 中，
\[
  P_{\text{naive}}(A)=\frac59,\quad
  P_{\text{naive}}(B)=\frac49,\quad
  P_{\text{naive}}(A\cup B)=\frac89,\quad
  P_{\text{naive}}(A\cap B)=\frac19.
\]
对于刚才考虑的这些事件的补集，
\[
  P_{\text{naive}}(A^c)=\frac49,\quad
  P_{\text{naive}}(B^c)=\frac59,\quad
  P_{\text{naive}}((A\cup B)^c)=\frac19,\quad
  P_{\text{naive}}((A\cap B)^c)=\frac89.
\]
一般地，
\[
  P_{\text{naive}}(A^c)
  =\frac{|A^c|}{|S|}
  =\frac{|S|-|A|}{|S|}
  =1-\frac{|A|}{|S|}
  =1-P_{\text{naive}}(A).
\]
在 1.6 节中我们会看到，关于补集的这个结果对概率总是成立，即使已经超出朴素定义的范围。求一个事件的概率时，一个好策略是先思考直接求该事件的概率更容易，还是求其补事件的概率更容易。De Morgan 律在这里尤其有用，因为有时处理交比处理并容易，反之亦然。

朴素定义非常受限，因为它要求 $S$ 有限，并且每颗鹅卵石具有相同质量。它常被误用：人们会在没有理由的情况下假设结果等可能，并作出“要么发生，要么不发生；我们不知道是哪一个，所以概率是 $50$--$50$”这类论证。除有时会给出荒谬概率之外，这种推理本身甚至也不自洽。例如，它会说火星上有生命的概率是 $1/2$（“要么有，要么没有”），同时也会说火星上有智慧生命的概率是 $1/2$；但凭直觉，以及 1.6 节中发展的概率性质，后者显然应当严格小于前者。不过，朴素定义确实适用于几类重要问题：
\begin{itemize}
  \item 当问题中有\textbf{对称性}，使各结果等可能时。例如，由于硬币的物理对称性，通常假设一枚硬币被抛掷后正面朝上的机会为 $50\%$。对于一副标准且充分洗匀的牌，假设所有排列等可能也是合理的；没有哪张特别积极的牌特别喜欢靠近牌堆顶端，牌堆中的任意位置都同样可能容纳 52 张牌中的任一张。
  \item 当结果是\textbf{按设计等可能}时。例如，考虑在一个有 $N$ 个人的总体中调查 $n$ 个人。一个常见目标是得到一个简单随机样本，也就是说，这 $n$ 个人是随机选择的，并且所有大小为 $n$ 的子集等可能。若能做到这一点，就保证朴素定义可用；但实践中这可能很难完成，因为存在各种复杂因素，例如无法取得总体中每个人完整而准确的联系信息。
  \item 当朴素定义作为一个有用的\textbf{零模型}时。在这种设定中，我们先假设朴素定义适用，看看它会给出什么预测，然后把观测数据与预测值比较，以评估“结果等可能”这一假设是否站得住脚。
\end{itemize}

\section*{1.4\quad 如何计数}

计算事件 $A$ 的朴素概率，需要数出 $A$ 中的鹅卵石数以及样本空间 $S$ 中的鹅卵石数。我们需要计数的集合常常极大。本节介绍一些基本计数方法；更进一步的方法可在组合数学中找到，组合数学就是研究计数的数学分支。

\subsection*{1.4.1\quad 乘法法则}

在一些问题中，我们可以用一个基本而灵活的原则直接数出可能性个数，这个原则称为\textbf{乘法法则}。我们会看到，乘法法则会自然导向有放回抽样和无放回抽样的计数规则，而这两种情形在概率与统计中经常出现。

\noindent\textbf{定理 1.4.1（乘法法则）}\quad 考虑一个由两个子实验组成的复合实验：实验 A 和实验 B。设实验 A 有 $a$ 个可能结果，并且对实验 A 的每一个结果，实验 B 都有 $b$ 个可能结果。那么这个复合实验共有 $ab$ 个可能结果。

为了看出乘法法则为什么成立，可以想象图 1.2 那样的树形图。让树按照实验 A 的可能性分出 $a$ 个分支，再对其中每一个分支按照实验 B 的可能性继续分出 $b$ 个分支。总共有
\[
  \underbrace{b+b+\cdots+b}_{a}=ab
\]
种可能。

\begin{figure}[htbp]
  \centering
  \includegraphics[width=0.42\textwidth]{1-2.png}
  \caption{说明乘法法则的树形图。如果实验 A 有 3 个可能结果，且对其中每一个结果，实验 B 有 4 个可能结果，那么总共有 $3\cdot 4=12$ 个可能结果。}
\end{figure}

\noindent\textbf{提醒 1.4.2}\quad 人们常常更容易按时间顺序思考实验，但乘法法则并不要求实验 A 必须先于实验 B 进行。

\noindent\textbf{例 1.4.3（赛跑）}\quad 假设有 10 个人参加赛跑。假定不会并列，并且 10 个人都会完成比赛，因此第一名、第二名和第三名都是明确的。第一、第二、第三名的获奖者有多少种可能？

\noindent\textbf{解}\quad 第一名有 10 种可能；一旦第一名确定，第二名有 9 种可能；一旦二者都确定，第三名有 8 种可能。因此由乘法法则，共有
\[
  10\cdot 9\cdot 8=720
\]
种可能。

我们不必先考虑第一名。也可以说，第三名有 10 种可能；一旦第三名确定，第二名有 9 种可能；一旦二者都确定，第一名有 8 种可能。或者想象有三个领奖台，赛后第一、第二和第三名分别站在金、银、铜三个领奖台上。仍然有 $10\cdot 9\cdot 8=720$ 种可能来填充这些领奖台，并没有理由必须按金、银、铜的顺序考虑领奖台。

\noindent\textbf{例 1.4.4（棋盘）}\quad 图 1.3 中的一个 $8\times 8$ 棋盘有多少个格子？仅“$8\times 8$ 棋盘”这个名称就使答案很明显：棋盘上有 $8\cdot 8=64$ 个格子。网格结构让这一点很清楚；也可以把它看作乘法法则的例子：要指定一个格子，可以指定它所在的行和列。行有 8 种选择，对每一行，列有 8 种选择。

进一步地，不做任何计算也能看出，白格和黑格各占一半。想象把棋盘顺时针旋转 $90^\circ$，则原来白格所在的位置现在都变成黑格，反之亦然，所以白格数必须等于黑格数。也可以用乘法法则数白格数：在 8 行中的每一行都有 4 个白格，因此共有 $8\cdot 4=32$ 个白格。

\begin{figure}[htbp]
  \centering
  \includegraphics[width=0.78\textwidth]{1-3.png}
  \caption{左边是一个 $8\times 8$ 棋盘，右边是一个填字游戏网格。棋盘有 $8\cdot 8=64$ 个格子，而要数出填字游戏网格中的白格数则需要更多工作。}
\end{figure}

相比之下，要数出图 1.3 所示填字游戏网格中的白格数需要更多努力。乘法法则在这里不适用，因为不同的行有时含有不同数量的白格。

\noindent\textbf{例 1.4.5（冰淇淋筒）}\quad 假设你要买一个冰淇淋筒。你可以选择蛋筒类型：普通蛋筒或华夫蛋筒；也可以选择口味：巧克力、香草或草莓。这个决策过程可以用图 1.4 中的树形图来表示。

\begin{figure}[htbp]
  \centering
  \includegraphics[width=0.76\textwidth]{1-4.png}
  \caption{选择冰淇淋筒的树形图。无论先选择蛋筒类型还是先选择口味，都共有 $2\cdot 3=3\cdot 2=6$ 种可能。}
\end{figure}

由乘法法则，共有 $2\cdot 3=6$ 种可能。这个例子很简单，但值得仔细思考，因为它是思考和可视化更复杂例子的基础。很快我们会遇到一些例子：如果真的把树画到可读大小，所需空间会超过已知宇宙的大小；但概念上我们仍然可以像冰淇淋例子那样思考。注意以下几点：
\begin{enumerate}
  \item 先选蛋筒类型（“我要一个装巧克力冰淇淋的华夫蛋筒”）还是先选口味（“我要一个放在华夫蛋筒上的巧克力冰淇淋”）并不重要。无论哪种方式，都有 $2\cdot 3=3\cdot 2=6$ 种可能。
  \item 普通蛋筒和华夫蛋筒是否提供同样的口味并不重要；重要的是对每一种蛋筒选择，都恰好有 3 种口味选择。如果出于某种奇怪的原因，华夫蛋筒上不能放巧克力冰淇淋，并且除香草和草莓之外没有替代口味，那么可能性就变成 $3+2=5$，乘法法则便不能适用。在更大的例子中，这类复杂性会使计数困难得多。
\end{enumerate}

现在假设某天你买了两个冰淇淋筒，一个在下午，一个在晚上。例如，用 $(\text{cakeC},\text{waffleV})$ 表示下午买了一个普通蛋筒巧克力冰淇淋，晚上买了一个华夫蛋筒香草冰淇淋。由乘法法则，这个美味的复合实验共有 $6^2=36$ 种可能。

但如果你只关心当天吃过哪些冰淇淋筒，而不关心吃的顺序，也就是说不想区分 $(\text{cakeC},\text{waffleV})$ 和 $(\text{waffleV},\text{cakeC})$，那现在是否有 $36/2=18$ 种可能？不是，因为像 $(\text{cakeC},\text{cakeC})$ 这样的可能原本就只列出一次。共有 $6\cdot 5=30$ 个有序可能 $(x,y)$ 满足 $x\ne y$，如果把 $(x,y)$ 与 $(y,x)$ 视为等价，它们变成 15 种可能；再加上 6 个形如 $(x,x)$ 的可能，共有 21 种可能。注意，如果原来的 36 个有序对 $(x,y)$ 等可能，那么这里的 21 种可能并不等可能。

\noindent\textbf{例 1.4.6（子集）}\quad 一个含 $n$ 个元素的集合有 $2^n$ 个子集，包括空集 $\varnothing$ 和它自身。这由乘法法则推出，因为对每个元素，我们都可以选择把它包含进子集中或排除在外。例如，集合 $\{1,2,3\}$ 有 8 个子集：
\[
  \varnothing,\ \{1\},\ \{2\},\ \{3\},\ \{1,2\},\ \{1,3\},\ \{2,3\},\ \{1,2,3\}.
\]
这个结果解释了为什么在例 1.2.3 中可以定义 $2^{52}$ 个事件。

我们可以用乘法法则推出有放回抽样和无放回抽样的公式。概率与统计中的许多实验都可以解释为这两种语境之一；因此，两条公式都直接来自同一个基本计数原则，是很自然的。

\noindent\textbf{定理 1.4.7（有放回抽样）}\quad 考虑 $n$ 个对象，并从中一次一个地作 $k$ 次选择，且有放回（即选中某个对象不会阻止它再次被选中）。那么共有 $n^k$ 个可能结果；这里顺序是重要的，例如先选对象 3 再选对象 7，与先选对象 7 再选对象 3，计为不同结果。

例如，想象一个罐子中有 $n$ 个球，标号为 $1$ 到 $n$。我们一次抽一个球并放回，也就是说每次抽中某个球后都把它放回罐中。每次抽球都是一个有 $n$ 个可能结果的子实验，共有 $k$ 个子实验。因此，由乘法法则，得到一个大小为 $k$ 的样本共有 $n^k$ 种方式。

\noindent\textbf{定理 1.4.8（无放回抽样）}\quad 考虑 $n$ 个对象，并从中一次一个地作 $k$ 次选择，且无放回（即选中某个对象会使它不能再次被选中）。若 $1\le k\le n$，则共有
\[
  n(n-1)\cdots(n-k+1)
\]
个可能结果；若 $k>n$，则有 0 个可能结果。这里顺序是重要的。按照约定，当 $k=1$ 时，$n(n-1)\cdots(n-k+1)=n$。

这个结果也直接来自乘法法则：每次抽到的球仍是一个子实验，而可能结果的数量每次减少 1。注意，从 $n$ 个对象中无放回抽取 $k$ 个时，需要 $k\le n$；而有放回抽样时，对象是取之不尽的。

\noindent\textbf{例 1.4.9（排列与阶乘）}\quad $1,2,\ldots,n$ 的一个\textbf{排列}，是把它们按某种顺序排成一列；例如，$3,5,1,2,4$ 是 $1,2,3,4,5$ 的一个排列。由定理 1.4.8 取 $k=n$，$1,2,\ldots,n$ 共有 $n!$ 个排列。例如，$n$ 个人排队买冰淇淋共有 $n!$ 种方式。回忆：对任意正整数 $n$，
\[
  n!=n(n-1)(n-2)\cdots 1,
\]
并且 $0!=1$。

定理 1.4.7 和定理 1.4.8 是关于\textbf{计数}的定理；但当朴素定义适用时，我们可以用它们来计算\textbf{概率}。这引出下一个例子：概率论中著名的生日问题。它的解法同时包含有放回抽样和无放回抽样。

\noindent\textbf{例 1.4.10（生日问题）}\quad 房间里有 $k$ 个人。假设每个人的生日等可能地落在一年 365 天中的任一天（排除 2 月 29 日），且不同人的生日相互独立（我们之后会正式定义独立性；直观地说，这意味着知道一些人的生日不会提供关于其他人生日的信息。例如，如果知道其中两人是双胞胎，这一点就不成立）。问：这群人中至少有一对人同一天生日的概率是多少？

\noindent\textbf{解}\quad 给房间里的人分配生日共有 $365^k$ 种方式，因为可以把一年 365 天想象为有放回地抽取 $k$ 次。根据假设，所有这些可能等可能，所以概率的朴素定义适用。

若直接使用朴素定义，我们需要数出给 $k$ 个人分配生日且至少两人同一天生日的方式数。但这个计数问题很难，因为同一天生日的可能是 Emma 和 Steve，也可能是 Steve 和 Naomi，也可能三人同一天生日，还可能三人同一天生日的同时，组中另外两人同在另一个生日，等等。

因此改为数补事件：给 $k$ 个人分配生日，使得没有两个人同一天生日的方式数。这相当于从一年 365 天中无放回地抽样，因此当 $k\le 365$ 时，可能性数量为
\[
  365\cdot 364\cdot 363\cdots(365-k+1).
\]
所以，在 $k$ 人组中没有生日相同的概率是
\[
  P(\text{没有生日相同})
  =\frac{365\cdot 364\cdots(365-k+1)}{365^k},
\]
而至少有一对生日相同的概率是
\[
  P(\text{至少一对生日相同})
  =1-\frac{365\cdot 364\cdots(365-k+1)}{365^k}.
\]

\begin{figure}[htbp]
  \centering
  \includegraphics[width=0.66\textwidth]{1-5.png}
  \caption{房间中有 $k$ 个人时，至少两人同一天生日的概率。该概率第一次超过 $0.5$ 发生在 $k=23$。}
\end{figure}

图 1.5 将至少有一对生日相同的概率画成 $k$ 的函数。使该概率超过 $0.5$ 的第一个 $k$ 值是 $k=23$。因此，在 23 个人组成的群体中，至少存在一对同生日者的机会超过 $50\%$。当 $k=57$ 时，同生日的概率已经超过 $99\%$。

当然，当 $k=366$ 时必然有一对生日相同；但令人惊讶的是，即使人数小得多，存在生日匹配的概率也已经压倒性地大。一个快速直觉是，23 个人中共有
\[
  \binom{23}{2}=253
\]
对人，其中任意一对都可能生日相同。

习题 26 和 27 说明，生日问题远不只是一个有趣的聚会游戏，也远不只是建立巧合直觉的方式；它在统计学和计算机科学中也有重要应用。习题 62 探讨更一般的情形，其中每一天的概率不一定都是 $1/365$。事实表明，在非等概率情形下，至少出现一次匹配甚至会更可能。

\noindent\textbf{提醒 1.4.11（给对象贴标签）}\quad 从总体中抽取样本是统计学中的一个非常基本的概念。重要的是，应当把总体中的对象或人看作有名字或有标签。例如，如果罐中有 $n$ 个球，我们可以想象它们带有从 $1$ 到 $n$ 的标签，即使在人眼看来这些球一模一样。在生日问题中，我们可以给每个人一个 ID 编号，而不是把人看作不可区分的粒子或一群没有面孔的人。

一个相关例子，是 Leibniz 在一个看似简单的问题中犯下的有启发性的错误。

\noindent\textbf{例 1.4.12（Leibniz 的错误）}\quad 如果掷两枚公平骰子，哪一个更可能：点数和为 11，还是点数和为 12？

\noindent\textbf{解}\quad 给两枚骰子贴上 A 与 B 的标签，并把每枚骰子视作一个子实验。由乘法法则，形如（A 的点数，B 的点数）的有序对共有 36 个可能结果，并且由对称性它们等可能。其中，$(5,6)$ 和 $(6,5)$ 有利于点数和为 11，而只有 $(6,6)$ 有利于点数和为 12。因此，和为 11 的可能性是和为 12 的两倍；前者概率为 $1/18$，后者概率为 $1/36$。

然而，Leibniz 错误地论证说，点数和为 11 与点数和为 12 等可能，因为这两个点数和都只能用一种方式得到。这里 Leibniz 犯的错误，是把两枚骰子当作不可区分的对象，将 $(5,6)$ 和 $(6,5)$ 看成同一个结果。

如何避免 Leibniz 的错误？第一，如提醒 1.4.11 所解释，我们应当给所考虑的对象贴标签，而不是把它们当作不可区分的对象。如果 Leibniz 给骰子贴上 A 和 B、绿色和橙色、左和右这样的标签，他就不会犯这个错误。第二，在用计数计算概率之前，应当先问朴素定义是否适用。

\subsection*{1.4.2\quad 校正过度计数}

在许多计数问题中，直接把每种可能恰好数一次并不容易。不过，如果我们能够把每种可能恰好数 $c$ 次，就可以通过除以 $c$ 来校正。例如，如果每种可能都恰好被重复数了两次，那么除以 2 就能得到正确计数。我们称这为\textbf{校正过度计数}。

\noindent\textbf{例 1.4.13（委员会与分队）}\quad 考虑一个由 4 人组成的群体。
\begin{enumerate}
  \item[(a)] 选出一个两人委员会有多少种方式？
  \item[(b)] 把 4 人分成两个两人队伍有多少种方式？
\end{enumerate}

\noindent\textbf{解}\quad
(a) 一种计数方式是直接列出所有可能。把 4 人标为 $1,2,3,4$，可能性为
\[
  \fbox{12},\ \fbox{13},\ \fbox{14},\ \fbox{23},\ \fbox{24},\ \fbox{34}.
\]
另一种方法是使用乘法法则并校正过度计数。由乘法法则，选择委员会第一人有 4 种方式，选择第二人有 3 种方式；但这把每种可能都数了两次，因为选 $1$ 和 $2$ 进入委员会，与选 $2$ 和 $1$ 进入委员会是同一件事。既然过度计数因子为 2，可能性数量就是 $(4\cdot 3)/2=6$。

(b) 有三种方式可以看出共有 3 种分队方法。仍把人标为 $1,2,3,4$，可以直接列出：
\[
  \fbox{12}\ \fbox{34},\qquad
  \fbox{13}\ \fbox{24},\qquad
  \fbox{14}\ \fbox{23}.
\]
不过人数更多时，直接列出会很快变得枯燥甚至不可行。另一种方法是注意到，只需指定 1 号的队友，另一个队伍就随之确定。第三种方法是利用 (a)：选出一个队伍有 6 种方式。这又把每种分队数了两次，因为选 $1$ 和 $2$ 成为一队，等价于选 $3$ 和 $4$ 成为一队。所以答案仍为 $6/2=3$。

\textbf{二项系数}计数的是某个集合中给定大小的子集个数，例如从 $n$ 个人中选出大小为 $k$ 的委员会的方式数。集合和子集按定义是无序的，例如 $\{3,1,4\}=\{4,1,3\}$；因此这里数的是从 $n$ 个对象中无放回选出 $k$ 个对象，且不区分这些对象被选出的不同顺序的方式数。

\noindent\textbf{定义 1.4.14（二项系数）}\quad 对任意非负整数 $k,n$，二项系数 $\binom{n}{k}$（读作“$n$ 选 $k$”）是一个大小为 $n$ 的集合中大小为 $k$ 的子集个数。

例如，如例 1.4.13 所示，$\binom42=6$。二项系数有时称为组合数；在这里我们不使用“组合”这个术语，因为“组合”是一个很有用的一般性词语。从代数上看，二项系数可以如下计算。

\noindent\textbf{定理 1.4.15（二项系数公式）}\quad 当 $k\le n$ 时，有
\[
  \binom{n}{k}
  =\frac{n(n-1)\cdots(n-k+1)}{k!}
  =\frac{n!}{(n-k)!k!}.
\]
当 $k>n$ 时，$\binom{n}{k}=0$。

\noindent\textbf{证明}\quad 设 $A$ 是一个满足 $|A|=n$ 的集合。$A$ 的任一子集大小至多为 $n$，所以当 $k>n$ 时 $\binom{n}{k}=0$。现在设 $k\le n$。由定理 1.4.8，无放回有序选择 $k$ 个元素有 $n(n-1)\cdots(n-k+1)$ 种方式。这把我们关心的每个子集都过度计数了 $k!$ 次（因为我们不关心这些元素的排列顺序），所以除以 $k!$ 得到正确计数。

\noindent\textbf{提醒 1.4.16}\quad 二项系数 $\binom{n}{k}$ 常用阶乘来定义，但要记住，即使负数的阶乘没有定义，当 $k>n$ 时仍有 $\binom{n}{k}=0$。此外，定理 1.4.15 中间的表达式往往比阶乘表达式更适合计算，因为阶乘增长极快。例如
\[
  \binom{100}{2}=\frac{100\cdot 99}{2}=4950
\]
甚至可以手算；而先计算 $100!$ 和 $98!$ 再算 $\binom{100}{2}=100!/(98!2!)$，则很浪费，且由于涉及极大的数（$100!\approx 9.33\times 10^{157}$）而可能很危险。

\noindent\textbf{例 1.4.17（俱乐部职员）}\quad 一个有 $n$ 人的俱乐部中，选出一名主席、一名副主席和一名财务主管共有 $n(n-1)(n-2)$ 种方式；而不预先指定头衔地选出 3 名职员共有
\[
  \binom{n}{3}=\frac{n(n-1)(n-2)}{3!}
\]
种方式。

\noindent\textbf{例 1.4.18（单词的排列）}\quad 单词 LALALAAA 的字母有多少种排列方式？要确定一个排列，只需选择 5 个 A 的位置（或等价地，决定 3 个 L 的位置）。因此共有
\[
  \binom85=\binom83=\frac{8\cdot 7\cdot 6}{3!}=56
\]
种排列。

单词 STATISTICS 的字母有多少种排列方式？有两种方法。可以先选择 S 的位置，再从剩余位置中选择 T 的位置，然后选择 I 的位置，再选择 A 的位置（此时 C 的位置也确定了）。或者，先从 $10!$ 开始，再校正过度计数：除以 $3!3!2!$，以考虑 S 之间可以任意互换，T 和 I 也同样可以互换。因此共有
\[
  \binom{10}{3}\binom73\binom42\binom21
  =\frac{10!}{3!3!2!}=50400
\]
种可能。

\noindent\textbf{例 1.4.19（二项式定理）}\quad 二项式定理说明，对任意非负整数 $n$，
\[
  (x+y)^n=\sum_{k=0}^n \binom{n}{k}x^k y^{n-k}.
\]
为了证明它，展开乘积
\[
  \underbrace{(x+y)(x+y)\cdots(x+y)}_{n\text{ 个因子}}.
\]
正如 $(a+b)(c+d)=ac+ad+bc+bd$ 是一些项的和，其中每一项从第一个因子中选 $a$ 或 $b$（但不同时选），从第二个因子中选 $c$ 或 $d$（但不同时选）；$(x+y)^n$ 的各项由从每个因子中选 $x$ 或 $y$（但不同时选）得到。恰好选择 $k$ 个 $x$ 有 $\binom{n}{k}$ 种方式，每一种选择给出项 $x^k y^{n-k}$。二项式定理随之成立。

在许多朴素定义适用的问题中，我们可以用二项系数计算概率。

\noindent\textbf{例 1.4.20（扑克中的葫芦）}\quad 从一副标准、充分洗匀的 52 张牌中发 5 张牌。如果一手牌由某个点数的三张牌和另一个点数的两张牌组成（例如三张 7 和两张 10，顺序任意），则称这手牌为扑克中的\textbf{葫芦}。抽到葫芦的概率是多少？

\noindent\textbf{解}\quad 所有 $\binom{52}{5}$ 种可能手牌由对称性等可能，所以朴素定义适用。为了求葫芦手牌的数量，使用乘法法则（并想象树形图）：先有 13 种选择来决定三张牌的点数；为具体起见，假设我们有三张 7，并关注树的这个分支。有 $\binom43$ 种方式选择是哪三张 7。然后有 12 种选择来决定两张牌的点数，例如具体地说是两张 10，并有 $\binom42$ 种方式选择是哪两张 10。因此
\[
  P(\text{葫芦})
  =\frac{13\binom43\,12\binom42}{\binom{52}{5}}
  =\frac{3744}{2598960}\approx 0.00144.
\]
玩扑克时，小数近似更有用；但用二项系数表达的答案是精确的，而且自带解释性（看到 $\binom{52}{5}$ 比看到 2598960 更能提示它的来源）。

\noindent\textbf{例 1.4.21（Newton--Pepys 问题）}\quad Samuel Pepys 为赌博目的向 Isaac Newton 咨询了下面的问题。下列事件中哪一个概率最大？
\[
\begin{aligned}
A&:\text{ 掷 6 枚公平骰子时，至少出现一个 6；}\\
B&:\text{ 掷 12 枚公平骰子时，至少出现两个 6；}\\
C&:\text{ 掷 18 枚公平骰子时，至少出现三个 6。}
\end{aligned}
\]

\noindent\textbf{解}\quad 三个实验分别有 $6^6,6^{12},6^{18}$ 个可能结果，并且由对称性，朴素定义在三个实验中都适用。

对 $A$，与其数出至少出现一个 6 的方式数，不如数出没有 6 的方式数。没有 6 等价于从数字 $1$ 到 $5$ 中有放回抽取 6 次，所以有 $5^6$ 个结果有利于 $A^c$（于是 $6^6-5^6$ 个结果有利于 $A$）。因此
\[
  P(A)=1-\frac{5^6}{6^6}\approx 0.67.
\]

对 $B$，同样先数 $B^c$ 中的结果。12 次掷骰没有 6 有 $5^{12}$ 种方式。恰好出现一个 6 有 $\binom{12}{1}5^{11}$ 种方式：先选哪一枚骰子为 6，再让其他 11 枚骰子从数字 $1$ 到 $5$ 中有放回抽样。相加得到未能至少出现两个 6 的方式数。因此
\[
  P(B)=1-\frac{5^{12}+\binom{12}{1}5^{11}}{6^{12}}\approx 0.62.
\]

对 $C$，数 $C^c$ 中的结果，也就是 18 次掷骰中出现零个、一个或两个 6 的方式数。没有 6 有 $5^{18}$ 种方式；恰好一个 6 有 $\binom{18}{1}5^{17}$ 种方式；恰好两个 6 有 $\binom{18}{2}5^{16}$ 种方式（选择哪两枚骰子为 6，再决定其他 16 枚骰子的点数）。于是
\[
  P(C)=1-\frac{5^{18}+\binom{18}{1}5^{17}+\binom{18}{2}5^{16}}{6^{18}}
  \approx 0.60.
\]
因此 $A$ 的概率最大。

Newton 用类似计算得到了正确答案。他还给 Pepys 提供了一个直觉论证来说明为什么 $A$ 是三者中最可能的；然而，这个直觉是无效的。使用有偏骰子可能导致 $A,B,C$ 的概率顺序不同，但 Newton 的直觉论证并不依赖骰子是否公平。

本书关心计数，并不是为了计数本身，而是因为计数有时能帮助我们求概率。下面是一个漂亮但危险的计数问题：解法很优雅，但其结果很少能和概率的朴素定义一起使用。

\noindent\textbf{例 1.4.22（Bose--Einstein）}\quad 如果顺序不重要（我们只关心每个对象被选中了多少次，而不关心选取顺序），从 $n$ 个对象中有放回地选择 $k$ 次，有多少种方式？

\noindent\textbf{解}\quad 当顺序重要时，由乘法法则答案是 $n^k$，但这里的问题困难得多。我们通过求解一个同构问题来解决它，也就是同一个问题的另一种伪装。

先求把 $k$ 个不可区分粒子放入 $n$ 个可区分盒子中的方式数。也就是说，以任意方式交换粒子都不视为一个不同可能；重要的只是每个盒子中有多少个粒子。这个情形称为 Bose--Einstein 问题，因为物理学家 Satyendra Nath Bose 和 Albert Einstein 在 20 世纪 20 年代研究了不可区分粒子的相关问题，并用他们的思想成功预测了一种称为 Bose--Einstein 凝聚的奇异物质状态的存在。

任一配置都可以自然地编码为由竖线和黑点组成的序列，如图 1.6 所示。

\begin{figure}[htbp]
  \centering
  \includegraphics[width=0.64\textwidth]{1-6.png}
  \caption{Bose--Einstein 编码：把 $k=7$ 个不可区分粒子放入 $n=4$ 个可区分盒子，可以表示为竖线和黑点组成的序列，其中竖线表示墙，黑点表示粒子。}
\end{figure}

一个有效序列必须以竖线开始和结束，并且在起始与结束竖线之间恰有 $n-1$ 条竖线和 $k$ 个黑点；反过来，任意这样的序列都有效地编码了某个粒子入盒配置。想象我们已经写下起始和结束竖线，它们表示外墙；在两道外墙之间，有 $n+k-1$ 个待填空位。只需选择其中哪些位置放入 $k$ 个黑点（于是内部 $n-1$ 条竖线的位置也完全确定）。所以可能性数量为
\[
  \binom{n+k-1}{k}.
\]
这种计数方法有时称为\textbf{星与棒}论证；在这里我们用黑点代替星号。

为了把这个结果联系回原问题，可以让每个盒子对应 $n$ 个对象中的一个，并用粒子作为“记号”来记录每个对象被选中了多少次。例如，若某个盒子恰好含有 3 个粒子，则意味着该盒子对应的对象被恰好选中了 3 次。粒子不可区分，对应的事实是：我们不关心对象被选中的顺序。因此，原问题的答案也是 $\binom{n+k-1}{k}$。

另一个同构问题，是数方程
\[
  x_1+x_2+\cdots+x_n=k
\]
的解 $(x_1,\ldots,x_n)$ 的个数，其中 $x_i$ 是非负整数。这是等价的，因为可以把 $x_i$ 看成第 $i$ 个盒子中的粒子数。

\noindent\textbf{提醒 1.4.23}\quad 除非在非常特殊的情形下，Bose--Einstein 结果不应与概率的朴素定义一起使用。例如，考虑一个调查：从大小为 $n$ 的总体中一次一个地、有放回地、以相等概率选择人，收集大小为 $k$ 的样本。那么 $n^k$ 个有序样本等可能，使朴素定义适用；但 $\binom{n+k-1}{k}$ 个无序样本（只关心每个人被抽中了多少次）并不等可能。

再举一个例子，令 $n=365$ 表示一年中的天数，$k$ 表示人数。无序的生日列表有多少种？例如当 $k=3$ 时，我们想数像（5 月 1 日，3 月 31 日，4 月 11 日）这样的列表，其中所有排列都视为等价。不能简单地用 $n^k/3!$ 之类的方法校正过度计数，因为（5 月 1 日，3 月 31 日，4 月 11 日）有 6 个排列，而（3 月 31 日，3 月 31 日，4 月 11 日）只有 3 个排列。由 Bose--Einstein，列表数为 $\binom{n+k-1}{k}$。但等可能的是有序生日列表，而不是无序生日列表，所以 Bose--Einstein 数值不应被用于计算生日概率。

\section*{1.5\quad 故事证明}

\textbf{故事证明}是通过解释来作出的证明。对计数问题来说，这通常意味着用两种不同方式数同一件事，而不是做繁琐代数。故事证明常常能避免混乱的计算，并且比代数证明更进一步，解释为什么结论为真。“故事”这个词有多种含义，有些比另一些更数学化；但这里所说的故事证明，是完全有效的数学证明。下面给出一些故事证明的例子，它们也同时是更多计数例子。

\noindent\textbf{例 1.5.1（选择补集）}\quad 对任意非负整数 $n,k$，若 $k\le n$，则
\[
  \binom{n}{k}=\binom{n}{n-k}.
\]
这很容易通过代数方式检查（把二项系数写成阶乘形式即可），但故事证明使结论更容易被直观理解。

\noindent\textbf{故事证明}\quad 考虑在 $n$ 人中选择一个大小为 $k$ 的委员会。我们知道共有 $\binom{n}{k}$ 种可能。另一种选择委员会的方式，是指定哪 $n-k$ 个人\textbf{不}在委员会中；指定谁在委员会中就决定了谁不在委员会中，反之亦然。因此等式两边相等，因为它们是对同一件事的两种计数方式。

\noindent\textbf{例 1.5.2（队长）}\quad 对任意正整数 $n,k$，若 $k\le n$，则
\[
  n\binom{n-1}{k-1}=k\binom{n}{k}.
\]
这同样很容易通过代数检查（利用任意正整数 $m$ 满足 $m!=m(m-1)!$），但故事证明更有洞察力。

\noindent\textbf{故事证明}\quad 考虑一个由 $n$ 人组成的群体，要从中选出一个 $k$ 人队伍，其中一人为队长。为了指定一种可能，可以先选队长，再选择剩余的 $k-1$ 名队员；这给出等式左边。等价地，也可以先选择 $k$ 名队员，再从中选择一人作为队长；这给出等式右边。

\noindent\textbf{例 1.5.3（Vandermonde 恒等式）}\quad 二项系数之间有一个著名关系，称为\textbf{Vandermonde 恒等式}：
\[
  \binom{m+n}{k}
  =\sum_{j=0}^{k}\binom{m}{j}\binom{n}{k-j}.
\]
这个恒等式在本书中会多次出现。如果试图用蛮力展开所有二项系数来证明它，会非常痛苦；但故事证明能优雅地证明结果，并清楚说明为什么恒等式成立。

\noindent\textbf{故事证明}\quad 考虑一个学生组织，其中有 $m$ 名低年级学生和 $n$ 名高年级学生，要从中选出一个大小为 $k$ 的委员会。共有 $\binom{m+n}{k}$ 种可能。如果委员会中有 $j$ 名低年级学生，那么委员会中必须有 $k-j$ 名高年级学生。恒等式右边正是在对所有可能的 $j$ 情形求和。

\noindent\textbf{例 1.5.4（配对）}\quad 我们用故事证明说明
\[
  \frac{(2n)!}{2^n\cdot n!}=(2n-1)(2n-3)\cdots 3\cdot 1.
\]

\noindent\textbf{故事证明}\quad 我们将说明两边都在数把 $2n$ 个人分成 $n$ 对的方式数。取 $2n$ 个人，并给他们编号 $1$ 到 $2n$。可以先把这些人排成某个顺序，然后说前两人是一对，接下来两人是一对，依此类推。这把每种配对过度计数了 $n!\cdot 2^n$ 次，因为各对之间的顺序不重要，每对内部的顺序也不重要。另一种方法是这样数可能性：1 号的搭档有 $2n-1$ 种选择，然后 2 号的搭档有 $2n-3$ 种选择（如果 2 号已经与 1 号配对，则改看 3 号），依此类推。

\section*{1.6\quad 概率的非朴素定义}

我们已经看到几种对样本空间中结果进行计数的方法；当朴素定义适用时，这些方法使我们能够计算概率。但朴素定义能带我们走的路有限，因为它要求结果等可能，也不能处理无限样本空间。为了推广概率概念，我们先写下一份简短的愿望清单，说明希望概率如何表现（在数学中，愿望清单上的条目称为公理），然后把概率函数定义为满足这些性质的函数。

下面是本书后续使用的概率一般定义。它只需要两个公理，但从这些公理出发可以证明大量关于概率的结果。

\noindent\textbf{定义 1.6.1（概率的一般定义）}\quad 一个\textbf{概率空间}由样本空间 $S$ 和概率函数 $P$ 组成。概率函数 $P$ 以事件 $A\subseteq S$ 为输入，并输出 $P(A)$，即一个介于 0 和 1 之间的实数。函数 $P$ 必须满足以下公理：
\begin{enumerate}
  \item $P(\varnothing)=0,\ P(S)=1$。
  \item 若 $A_1,A_2,\ldots$ 是两两不交事件，则
  \[
    P\left(\bigcup_{j=1}^{\infty}A_j\right)
    =\sum_{j=1}^{\infty}P(A_j).
  \]
\end{enumerate}
这里说这些事件两两不交，是指它们互斥：当 $i\ne j$ 时，$A_i\cap A_j=\varnothing$。

在鹅卵石世界中，这个定义说概率像质量一样表现：空鹅卵石堆的质量为 0，所有鹅卵石的总质量为 1；如果有一些互不重叠的鹅卵石堆，则它们的合质量可以通过把各堆质量相加得到。与朴素情形不同，现在鹅卵石可以有不同质量；也可以有可数无穷多颗鹅卵石，只要总质量为 1。

我们甚至可以有不可数样本空间，例如令 $S$ 是平面上的一个区域。在这种情形下，不再想象鹅卵石，而可以想象泥浆铺满某个区域，其中泥浆总质量为 1。

任意函数 $P$，只要它把事件映射到区间 $[0,1]$ 中的数，并满足上述两个公理，就被视为有效的概率函数。不过，公理并不告诉我们应当如何\textbf{解释}概率；对此存在不同学派。

概率的\textbf{频率学派}观点认为，概率表示一次实验经过大量重复后的长期频率：如果我们说一枚硬币正面朝上的概率是 $1/2$，这意味着如果反复抛掷这枚硬币，正面朝上的频率会是 $50\%$。

概率的\textbf{贝叶斯学派}观点认为，概率表示关于所讨论事件的信念程度，因此我们可以给“候选人 A 会赢得选举”或“被告有罪”这样的假设赋予概率，即使不可能一遍又一遍地重复同一次选举或同一起案件。

贝叶斯观点和频率观点是互补的，二者都有助于在后续章节中发展直觉。不论选择如何解释概率，我们都可以使用这两个公理推出概率的其他性质，而且这些结果对任意有效的概率函数都成立。

\noindent\textbf{定理 1.6.2（概率的性质）}\quad 对任意事件 $A$ 和 $B$，概率具有以下性质：
\begin{enumerate}
  \item $P(A^c)=1-P(A)$。
  \item 若 $A\subseteq B$，则 $P(A)\le P(B)$。
  \item $P(A\cup B)=P(A)+P(B)-P(A\cap B)$。
\end{enumerate}

\noindent\textbf{证明}\quad
1. 因为 $A$ 与 $A^c$ 两两不交，且并为 $S$，第二公理给出
\[
  P(S)=P(A\cup A^c)=P(A)+P(A^c).
\]
而由第一公理 $P(S)=1$，所以 $P(A)+P(A^c)=1$。

2. 若 $A\subseteq B$，则可以把 $B$ 写成 $A$ 与 $B\cap A^c$ 的并，其中 $B\cap A^c$ 是 $B$ 中不属于 $A$ 的部分。下图说明了这一点。

\begin{figure}[htbp]
  \centering
  \includegraphics[width=0.42\textwidth]{1-6-1.png}
\end{figure}

由于 $A$ 与 $B\cap A^c$ 两两不交，可以应用第二公理：
\[
  P(B)=P(A\cup(B\cap A^c))=P(A)+P(B\cap A^c).
\]
概率非负，所以 $P(B\cap A^c)\ge0$，从而 $P(B)\ge P(A)$。

3. 这个结果的直觉可以用如下 Venn 图看出。

\begin{figure}[htbp]
  \centering
  \includegraphics[width=0.44\textwidth]{1-6-2.png}
\end{figure}

阴影区域表示 $A\cup B$，但这个区域的概率不是 $P(A)+P(B)$，因为那会把橄榄球形状的区域 $A\cap B$ 数两次。为了校正这一点，我们减去 $P(A\cap B)$。这是有用的直觉，但还不是证明。

为了用概率公理证明，可以把 $A\cup B$ 写成两个不交事件 $A$ 与 $B\cap A^c$ 的并。于是由第二公理，
\[
  P(A\cup B)=P(A\cup(B\cap A^c))=P(A)+P(B\cap A^c).
\]
所以只需说明 $P(B\cap A^c)=P(B)-P(A\cap B)$。由于 $A\cap B$ 和 $B\cap A^c$ 两两不交，且它们的并为 $B$，再应用第二公理得到
\[
  P(A\cap B)+P(B\cap A^c)=P(B).
\]
因此 $P(B\cap A^c)=P(B)-P(A\cap B)$，证毕。

第三条性质是\textbf{容斥原理}的一个特殊情形。容斥原理是在事件不一定两两不交时求事件并概率的公式。上面已经说明，对两个事件 $A$ 和 $B$，
\[
  P(A\cup B)=P(A)+P(B)-P(A\cap B).
\]
对三个事件，容斥原理给出
\[
\begin{aligned}
  P(A\cup B\cup C)
  &=P(A)+P(B)+P(C)\\
  &\quad -P(A\cap B)-P(A\cap C)-P(B\cap C)\\
  &\quad +P(A\cap B\cap C).
\end{aligned}
\]
为了获得直觉，可以考虑如下三重 Venn 图。

\begin{figure}[htbp]
  \centering
  \includegraphics[width=0.42\textwidth]{1-6-3.png}
\end{figure}

为了得到阴影区域 $A\cup B\cup C$ 的总面积，我们先加上三个圆的面积 $P(A)+P(B)+P(C)$。三个橄榄球形状的区域各被数了两次，所以再减去 $P(A\cap B)+P(A\cap C)+P(B\cap C)$。最后，中心区域被加了三次又减了三次；为了恰好数一次，必须把它加回来。这样保证图中的每个区域都被恰好数一次。

现在可以写出 $n$ 个事件的容斥公式。

\noindent\textbf{定理 1.6.3（容斥原理）}\quad 对任意事件 $A_1,\ldots,A_n$，
\[
\begin{aligned}
  P\left(\bigcup_{i=1}^{n}A_i\right)
  &=\sum_i P(A_i)-\sum_{i<j}P(A_i\cap A_j)
    +\sum_{i<j<k}P(A_i\cap A_j\cap A_k)-\cdots\\
  &\quad +(-1)^{n+1}P(A_1\cap\cdots\cap A_n).
\end{aligned}
\]
这个公式可以只用公理和归纳法证明；这里暂不展开。一般公式中加减交替的理由，与我们已经考虑过的特殊情形类似。

下一个例子，即 de Montmort 的匹配问题，是容斥原理的著名应用。Pierre Remond de Montmort 是一位法国数学家，他在赌博背景中研究概率，并写了一部专门分析各类纸牌游戏的著作。他在 1708 年提出了下面的问题，这个问题基于一种名为 Treize 的纸牌游戏。

\noindent\textbf{例 1.6.4（de Montmort 匹配问题）}\quad 考虑一副充分洗匀的 $n$ 张牌，标号为 $1$ 到 $n$。你一张一张翻开这些牌，同时依次说出数字 $1$ 到 $n$。如果在某一时刻，你说出的数字恰好等于正在翻开的牌上的数字（例如牌堆中的第 7 张牌标号为 7），就赢得游戏。获胜概率是多少？

\noindent\textbf{解}\quad 令 $A_i$ 为牌堆中第 $i$ 张牌标号为 $i$ 的事件。我们关心并集 $A_1\cup\cdots\cup A_n$ 的概率：只要至少有一张牌的号码与其在牌堆中的位置相同，就赢得游戏。使你输掉游戏的排列称为错排。

为了求并集概率，我们使用容斥原理。首先，对所有 $i$，
\[
  P(A_i)=\frac1n.
\]
一种看法是使用概率的朴素定义和完整样本空间：牌堆共有 $n!$ 个等可能排列，其中 $(n-1)!$ 个有利于 $A_i$（固定编号为 $i$ 的牌位于第 $i$ 个位置，其余 $n-1$ 张牌任意排列）。另一种看法是利用对称性：编号为 $i$ 的牌等可能地位于牌堆的 $n$ 个位置中的任意一个，因此它位于第 $i$ 个位置的概率为 $1/n$。其次，
\[
  P(A_i\cap A_j)=\frac{(n-2)!}{n!}=\frac{1}{n(n-1)},
\]
因为要求编号为 $i$ 和 $j$ 的牌分别位于第 $i$ 和第 $j$ 个位置。
牌堆中这两张牌固定后，其余 $n-2$ 张牌可以任意排列，所以有 $(n-2)!$ 个有利于 $A_i\cap A_j$ 的结果（总可能数为 $n!$）。类似地，
\[
  P(A_i\cap A_j\cap A_k)=\frac{1}{n(n-1)(n-2)},
\]
并且这个模式会继续适用于 4 个事件的交，等等。

在容斥公式中，涉及一个事件的项有 $n$ 个，涉及两个事件的项有 $\binom n2$ 个，涉及三个事件的项有 $\binom n3$ 个，依此类推。由问题的对称性，所有形如 $P(A_i)$ 的 $n$ 项都相等，所有形如 $P(A_i\cap A_j)$ 的 $\binom n2$ 项都相等，于是整个表达式大幅简化：
\[
\begin{aligned}
  P\left(\bigcup_{i=1}^{n}A_i\right)
  &=\frac nn-\frac{\binom n2}{n(n-1)}
    +\frac{\binom n3}{n(n-1)(n-2)}
    -\cdots+(-1)^{n+1}\frac1{n!}\\
  &=1-\frac1{2!}+\frac1{3!}-\cdots+(-1)^{n+1}\frac1{n!}.
\end{aligned}
\]
把它与 $1/e$ 的 Taylor 级数
\[
  e^{-1}=1-\frac1{1!}+\frac1{2!}-\frac1{3!}+\cdots
\]
比较可知，当 $n$ 很大时，获胜概率极接近 $1-1/e$，约为 $0.63$。有趣的是，随着 $n$ 增大，获胜概率趋近于 $1-1/e$，而不是趋近于 0 或 1。当牌堆中有很多牌时，可能匹配的位置数量增加，但任意特定匹配的概率下降，这两种力量相互抵消，平衡后得到约 $1-1/e$ 的概率。

容斥原理是求事件并概率的非常一般的公式，但当事件 $A_j$ 之间存在对称性时，它最有帮助；否则求和会极其繁琐。一般而言，在缺少对称性时，应当先尝试其他工具，把容斥原理作为最后手段。

\section*{1.7\quad 小结}

概率使我们能够有原则地量化不确定性和随机性。概率出现在我们进行实验时：实验所有可能结果的集合称为样本空间，而样本空间的子集称为事件。在自然语言中描述事件和把事件数学地写成集合（常用并、交和补）之间来回转换，是很有用的能力。

当样本空间有限时，鹅卵石世界能帮助我们可视化样本空间和事件。在鹅卵石世界中，每个结果是一颗鹅卵石，每个事件是一组鹅卵石。如果所有鹅卵石质量相同（即等可能），就可以应用概率的朴素定义，用计数来计算概率。

为此，我们讨论了若干计数工具。计数可能性数量时，我们常用乘法法则。例如，数字 $1,2,\ldots,n$ 有 $n!$ 个排列，含 $n$ 个元素的集合有 $2^n$ 个子集。如果不能直接使用乘法法则，有时可以把每种可能恰好数 $c$ 次，然后除以 $c$ 得到实际可能性数量。例如，这个策略有助于用阶乘表示二项系数。

一个重要陷阱，是误用概率的朴素定义：没有根据地显式或隐式假设结果等可能。一个有助于避免这一点的技巧，是给对象贴标签，以便表述精确，也避免把它们误当作不可区分的对象。

超越朴素定义后，我们把概率定义为一个函数，它以事件为输入，并给事件赋予一个介于 0 和 1 之间的实数。我们要求有效的概率函数满足两个公理：
\[
  P(\varnothing)=0,\qquad P(S)=1,
\]
以及：如果 $A_1,A_2,\ldots$ 两两不交，则
\[
  P\left(\bigcup_{j=1}^{\infty}A_j\right)
  =\sum_{j=1}^{\infty}P(A_j).
\]
许多有用性质都可以从这些公理推出。例如，对任意事件 $A$，
\[
  P(A^c)=1-P(A),
\]
并且有容斥公式
\[
\begin{aligned}
  P\left(\bigcup_{i=1}^{n}A_i\right)
  &=\sum_i P(A_i)-\sum_{i<j}P(A_i\cap A_j)
    +\sum_{i<j<k}P(A_i\cap A_j\cap A_k)-\cdots\\
  &\quad +(-1)^{n+1}P(A_1\cap\cdots\cap A_n)
\end{aligned}
\]
对任意事件 $A_1,\ldots,A_n$ 成立。当 $n=2$ 时，它变成更简洁的形式
\[
  P(A_1\cup A_2)=P(A_1)+P(A_2)-P(A_1\cap A_2).
\]

图 1.7 展示了概率函数如何把事件映射为 0 到 1 之间的数。随着我们继续概率之旅，还会向这张图中加入许多新概念。

\begin{figure}[htbp]
  \centering
  \includegraphics[width=0.82\textwidth]{1-7.png}
  \caption{区分事件和概率很重要。前者是集合，后者是数字。实验完成之前，我们通常不知道某个特定事件是否会发生，因此使用概率函数 $P$ 给它赋予发生概率。可以用集合运算根据旧事件定义新事件，也可以用概率性质把新事件的概率与旧事件的概率联系起来。}
\end{figure}

\section*{1.8\quad 习题}

标有 (S) 的习题在原书配套网站有详细解答。

\subsection*{计数}

\begin{enumerate}
  \item 单词 MISSISSIPPI 的字母有多少种排列方式？

  \item
  \begin{enumerate}
    \item 假设第一位不能是 0 或 1，问有多少个 7 位电话号码？
    \item 重新求解 (a)，但再假设电话号码也不允许以 911 开头。
  \end{enumerate}

  \item Fred 计划在某一周从周一到周五每天晚上出去吃饭，每顿饭都在他最喜欢的 10 家餐馆之一。
  \begin{enumerate}
    \item 如果 Fred 不愿在同一家餐馆吃超过一次，那么周一到周五的晚餐安排有多少种可能？
    \item 如果 Fred 愿意在同一家餐馆吃超过一次，但不愿连续两次（或更多次）在同一家餐馆吃，那么有多少种可能？
  \end{enumerate}

  \item 有 $n$ 名网球选手参加循环赛，也就是说每名选手都与其他每名选手恰好比赛一次。
  \begin{enumerate}
    \item 这个锦标赛有多少种可能结果（结果列出每一场谁赢谁输）？
    \item 总共进行了多少场比赛？
  \end{enumerate}

  \item 有 $2^n$ 名网球选手参加淘汰赛。每一轮中，胜者进入下一轮，败者被淘汰，直到只剩一人。例如，若起初有 $2^4=16$ 名选手，则第一轮有 8 场比赛，8 名胜者进入第二轮，再有 4 名胜者进入第三轮，再有 2 名胜者进入第四轮，最后一场的胜者成为冠军。
  \begin{enumerate}
    \item 共有多少轮？
    \item 通过把各轮比赛数相加，数出总共进行了多少场比赛。
    \item 这一次不做几乎任何计算，直接思考总共进行了多少场比赛。提示：需要淘汰多少名选手？
  \end{enumerate}

  \item 某天一个国际象棋俱乐部有 20 人。他们各自寻找对手并开始下棋。若在每局棋中谁执白无关紧要，问配对方式有多少种？

  \item 两名棋手 A 和 B 要下 7 局棋。每局有三种可能结果：A 胜、平局、A 负。A 胜得 1 分，平局得 $0.5$ 分，A 负得 0 分。
  \begin{enumerate}
    \item 若最终 A 有 3 胜 2 平 2 负，单局结果序列有多少种可能？
    \item 若最终 A 得 4 分、B 得 3 分，单局结果序列有多少种可能？
    \item 现在假设他们进行七局四胜制比赛：当某人达到 4 分或已经下完 7 局时，比赛结束，以先发生者为准。例如，若 6 局后 A 以 4 比 2 领先，则 A 获胜且不再下第 7 局。若比赛持续 7 局且 A 以 4 比 3 获胜，单局结果序列有多少种可能？
  \end{enumerate}

  \item (S)
  \begin{enumerate}
    \item 把 12 个人分成 3 个队，其中一个队有 2 人，另外两个队各有 5 人，有多少种方式？
    \item 把 12 个人分成 3 个队，每队 4 人，有多少种方式？
  \end{enumerate}

  \item (S)
  \begin{enumerate}
    \item 在平面中，从点 $(0,0)$ 到点 $(110,111)$ 有多少条路径，其中每一步都向上走一个单位或向右走一个单位？
    \item 从 $(0,0)$ 到 $(210,211)$ 有多少条路径，每一步都向上或向右走一个单位，并且路径必须经过 $(110,111)$？
  \end{enumerate}

  \item 为满足某个学位要求，学生可从 20 门课程中任选 7 门，但要求这 7 门中至少有 1 门必须是统计课。假设 20 门课中有 5 门统计课。
  \begin{enumerate}
    \item 选择 7 门课有多少种方式？
    \item 直观解释为什么 (a) 的答案不是 $\binom51\binom{19}{6}$。
  \end{enumerate}

  \item 设 $A$ 和 $B$ 是集合，$|A|=n,\ |B|=m$。
  \begin{enumerate}
    \item 从 $A$ 到 $B$ 的函数有多少个？
    \item 从 $A$ 到 $B$ 的一一函数有多少个？
  \end{enumerate}

  \item A、B、C、D 四名玩家玩纸牌游戏。一副标准且充分洗匀的牌发给四名玩家（每人 13 张）。
  \begin{enumerate}
    \item 玩家 A 拿到的手牌有多少种可能？（一手牌内部收牌顺序不重要。）
    \item 若关心哪名玩家拿到哪手牌，但不关心每手牌内部顺序，则所有人手牌的总体可能有多少种？
    \item 直观解释为什么 (b) 的答案不是 (a) 答案的四次方。
  \end{enumerate}

  \item 某赌场把 10 副标准牌混在一起成为一副大牌，称为 superdeck。因此这副大牌有 $52\cdot 10=520$ 张牌，每张牌有 10 个副本。从这副大牌中发出 10 张牌，不计顺序，也不区分牌来自原来的哪一副，问不同的 10 张手牌有多少种？请把答案表示为二项系数。提示：Bose--Einstein。

  \item 你要订两张披萨。披萨尺寸可以是小、中、大或特大，可以选择 8 种配料的任意组合（可以没有配料，也可以 8 种都有）。你的两张披萨共有多少种可能？
\end{enumerate}

\subsection*{故事证明}

\begin{enumerate}
  \setcounter{enumi}{14}
  \item (S) 给出故事证明：
  \[
    \sum_{k=0}^{n}\binom nk=2^n.
  \]

  \item (S) 证明对所有正整数 $n,k$ 且 $n\ge k$，
  \[
    \binom nk+\binom n{k-1}=\binom{n+1}{k},
  \]
  要分别用 (a) 代数方法和 (b) 故事证明，并解释为什么两边数的是同一件事。故事证明提示：想象一个有 $n+1$ 人的组织，其中一人被预先指定为主席。

  \item 给出故事证明：
  \[
    \sum_{k=0}^{n}\binom nk^2=\binom{2n}{n},
  \]
  对所有正整数 $n$ 成立。

  \item 给出故事证明：
  \[
    \sum_{k=1}^{n} k\binom nk^2
    =n\binom{2n-1}{n-1},
  \]
  对所有正整数 $n$ 成立。提示：考虑从两个大小均为 $n$ 的群体中选择一个大小为 $n$ 的委员会，其中只有其中一个群体的人有资格成为委员会主席。

  \item 给出故事证明：
  \[
    \sum_{k=2}^{n}\binom k2\binom{n-k+2}{2}
    =\binom{n+3}{5},
  \]
  对所有整数 $n\ge2$ 成立。提示：考虑集合 $\{1,2,\ldots,n+3\}$ 的大小为 5 的子集中间那个数。

  \item (S)
  \begin{enumerate}
    \item 用故事证明说明
    \[
      \binom kk+\binom{k+1}{k}+\binom{k+2}{k}+\cdots+\binom nk
      =\binom{n+1}{k+1},
    \]
    其中 $n,k$ 为正整数且 $n\ge k$。这称为\textbf{冰球杆恒等式}。提示：想象把一群人按年龄排序，然后考虑所选子群中年龄最大的人。
    \item 一大包 Haribo 小熊软糖可能有 30 到 50 颗。共有 5 种美味口味：菠萝（透明）、树莓（红）、橙子（橙）、草莓（绿色）、柠檬（黄）。没有不好吃的口味。这样的软糖包组成有多少种可能？答案可以用两个二项系数表示，但不要写成许多二项系数之和。
  \end{enumerate}

  \item 定义 $\left\{\begin{matrix} n\\ k\end{matrix}\right\}$ 为把 $\{1,2,\ldots,n\}$ 划分成 $k$ 个非空子集的方式数，也就是把 $n$ 名学生分成 $k$ 个组且每组至少一人的方式数。例如，$\left\{\begin{matrix}4\\2\end{matrix}\right\}=7$，因为有
  \[
  \begin{gathered}
    \{1\},\{2,3,4\};\quad \{2\},\{1,3,4\};\quad
    \{3\},\{1,2,4\};\quad \{4\},\{1,2,3\};\\
    \{1,2\},\{3,4\};\quad \{1,3\},\{2,4\};\quad
    \{1,4\},\{2,3\}.
  \end{gathered}
  \]
  证明恒等式：
  \[
    \left\{\begin{matrix} n+1\\ k\end{matrix}\right\}
    =
    \left\{\begin{matrix} n\\ k-1\end{matrix}\right\}
    +k\left\{\begin{matrix} n\\ k\end{matrix}\right\}.
  \]
  提示：我或者独自成组，或者不是。

  证明恒等式：
  \[
    \sum_{j=k}^{n}\binom nj
    \left\{\begin{matrix} j\\ k\end{matrix}\right\}
    =
    \left\{\begin{matrix} n+1\\ k+1\end{matrix}\right\}.
  \]
  提示：先决定有多少人不在我的组里。

  \item 荷兰数学家 R.J. Stroeker 评论说：每个数论初学者一定都会惊叹于这样一个奇妙事实：对每个自然数 $n$，前 $n$ 个正整数的立方和是一个完全平方数。而且，它正是前 $n$ 个正整数之和的平方：
  \[
    1^3+2^3+\cdots+n^3=(1+2+\cdots+n)^2.
  \]
  通常这个恒等式用归纳法证明，但那并不能给出太多解释。本题要求给出故事证明。
  \begin{enumerate}
    \item 给出恒等式
    \[
      1+2+\cdots+n=\binom{n+1}{2}
    \]
    的故事证明。提示：考虑循环赛。
    \item 给出恒等式
    \[
      1^3+2^3+\cdots+n^3
      =6\binom{n+1}{4}+6\binom{n+1}{3}+\binom{n+1}{2}
    \]
    的故事证明。提示：想象先在 $1$ 到 $n$ 之间选一个数，再有放回地从 $0$ 到比原数小的数中选 3 个数，然后按选出的不同数个数分类。
  \end{enumerate}
\end{enumerate}

\subsection*{概率的朴素定义}

\begin{enumerate}
  \setcounter{enumi}{22}
  \item 三个人在一栋 10 层楼的一楼进入空电梯。每个人按下自己想去楼层的按钮（除非别人已经按过）。假设他们等可能地想去 2 到 10 层，且彼此独立。问：按下的按钮对应 3 个连续楼层的概率是多少？

  \item (S) 某家庭有 6 个孩子，其中 3 个男孩、3 个女孩。假设所有出生顺序等可能，问年龄最大的 3 个孩子都是女孩的概率是多少？

  \item (S) 一个有 6 个区的城市在某周发生 6 起抢劫。假设抢劫发生地点随机，且所有“哪起抢劫发生在哪个区”的可能等可能。问：某个区发生超过 1 起抢劫的概率是多少？

  \item 某城市有 100 万居民。调查全部居民太昂贵，因此随机选择大小为 1000 的样本。本题假设调查是一次一个、有放回、等概率地选择人。
  \begin{enumerate}
    \item 解释这里的有放回抽样和无放回抽样与生日问题有什么关系。
    \item 求至少有一个人被选中超过一次的概率。
  \end{enumerate}

  \item 哈希表是计算机科学中常用的数据结构，能快速检索信息。设要存储一些人的电话号码，且没有两人同名。对每个名字 $x$，使用哈希函数 $h$，令 $h(x)$ 为存储 $x$ 电话号码的位置。假设有 $k$ 个人，每个人的电话号码被存储到 $1$ 到 $n$ 之间的一个随机位置，且各位置等概率并独立。求至少有一个位置存储超过一个电话号码的概率。

  \item (S) 某学院有 10 个课程时段，并把课程独立地完全随机分配到时段。学院恰好开设 3 门统计课。问：至少两门统计课在同一时段的概率是多少？

  \item (S) 对每一小题，判断空格中应填 $=,<,>$ 中哪一个，并给出清楚解释。
  \begin{enumerate}
    \item 掷 4 枚公平骰子总和为 21 的概率 \underline{\hspace{1cm}} 掷 4 枚公平骰子总和为 22 的概率。
    \item 随机 2 字母词是回文的概率 \underline{\hspace{1cm}} 随机 3 字母词是回文的概率。
  \end{enumerate}

  \item 沿用上一题定义，求随机 $n$ 字母词为回文的概率，分别取 $n=7$ 和 $n=8$。

  \item (S) 某森林中有 $N$ 只麋鹿，先捕获并标记一个大小为 $n$ 的简单随机样本，然后放回总体；之后再抽取一个大小为 $m$ 的新样本。这是生态学中常用的捕获--再捕获方法。求新样本中恰好有 $k$ 只先前被标记过的概率。

  \item 桌上有 4 张牌面朝下。你被告知其中两张红、两张黑，需要猜哪两张红、哪两张黑。你通过指认两张牌为红来猜测。假设所有配置等可能且你没有超能力，求恰好 $j$ 个猜测正确的概率，其中 $j=0,1,2,3,4$。

  \item (S) 罐中有 $r$ 个红球和 $g$ 个绿球，其中 $r,g$ 为固定正整数。随机抽出一个球，再随机抽出第二个球。
  \begin{enumerate}
    \item 直观解释为什么第二个球为绿球的概率与第一个球为绿球的概率相同。
    \item 为本问题的样本空间定义记号，用它计算 (a) 中的概率并说明它们相同。
    \item 假设总共有 16 个球，且两球同色的概率与两球异色的概率相同。求 $r$ 和 $g$ 的所有可能。
  \end{enumerate}

  \item (S) 从一副标准牌中随机发 5 张扑克手牌。用二项系数表示以下概率：
  \begin{enumerate}
    \item 同花（5 张牌花色相同；不计算皇家同花顺，即 A、K、Q、J、10 同花）。
    \item 两对（例如两张 3、两张 7 和一张 A）。
  \end{enumerate}

  \item 从一副标准牌中随机发 13 张手牌。求这手牌至少有 3 张来自每种花色的概率。

  \item 掷 30 枚骰子。求 $1,2,3,4,5,6$ 每个点数都恰好出现 5 次的概率。

  \item 一副牌充分洗匀，然后逐张发牌，直到第一次出现 A。
  \begin{enumerate}
    \item 求第一个 A 之前没有 K、Q、J 出现的概率。
    \item 求第一个 A 之前恰好出现一张 K、一张 Q、一张 J（顺序任意）的概率。
  \end{enumerate}

  \item Tyrion、Cersei 和另外 10 个人围坐在圆桌旁，座位随机安排。求 Tyrion 和 Cersei 相邻而坐的概率，用两种方法求：
  \begin{enumerate}
    \item 使用大小为 $12!$ 的样本空间，其中结果完整描述座位安排；
    \item 使用一个小得多的样本空间，只关注 Tyrion 和 Cersei。
  \end{enumerate}

  \item 一个组织有 $2n$ 人，由 $n$ 对夫妻组成。随机选出大小为 $k$ 的委员会，所有可能等可能。求委员会中恰好有 $j$ 对夫妻的概率。

  \item 罐中有 $n$ 个球，标号为 $1,2,\ldots,n$。有放回地逐个抽取共 $k$ 个球，得到一个数字序列。
  \begin{enumerate}
    \item 所得序列严格递增的概率是多少？
    \item 所得序列递增（但不一定严格递增，即允许重复）的概率是多少？
  \end{enumerate}

  \item 每个 $n$ 个球独立地放入 $n$ 个盒子之一，所有盒子等可能。求恰好一个盒子为空的概率。

  \item (S) 一个 norepeatword 是由通常的 26 个小写字母中的至少一个（也可能全部）组成的序列，且不允许重复。例如 course 是 norepeatword，而 statistics 不是。顺序重要，例如 course 与 source 不同。随机选择一个 norepeatword，且所有 norepeatword 等可能。证明它使用全部 26 个字母的概率非常接近 $1/e$。
\end{enumerate}

\subsection*{概率公理}

\begin{enumerate}
  \setcounter{enumi}{42}
  \item 证明对任意事件 $A,B$，
  \[
    P(A)+P(B)-1\le P(A\cap B)\le P(A\cup B)\le P(A)+P(B).
  \]
  对这三个不等式中的每一个，给出它成为等式的简单判据。

  \item 设 $A,B$ 为事件，差集 $B-A$ 定义为 $B$ 中所有不属于 $A$ 的元素。直接使用概率公理证明：若 $A\subseteq B$，则
  \[
    P(B-A)=P(B)-P(A).
  \]

  \item 设 $A,B$ 为事件。对称差 $A\triangle B$ 定义为属于 $A$ 或 $B$ 但不同时属于二者的元素集合。在逻辑和工程中，它也称为 $A$ 与 $B$ 的 XOR。直接使用概率公理证明：
  \[
    P(A\triangle B)=P(A)+P(B)-2P(A\cap B).
  \]

  \item 设 $A_1,A_2,\ldots,A_n$ 为事件。令 $B_k$ 为恰好 $k$ 个 $A_i$ 发生的事件，令 $C_k$ 为至少 $k$ 个 $A_i$ 发生的事件，$0\le k\le n$。用 $P(C_k)$ 和 $P(C_{k+1})$ 表示 $P(B_k)$。

  \item 事件 $A,B$ 独立，是指 $P(A\cap B)=P(A)P(B)$。
  \begin{enumerate}
    \item 在有限样本空间 $S$ 中给出两个独立事件 $A,B$ 的例子（两者都不等于 $\varnothing$ 或 $S$），并用鹅卵石世界图说明。
    \item 考虑在矩形
    \[
      R=\{(x,y):0<x<1,\ 0<y<1\}
    \]
    中随机选点，点落入 $R$ 内某个区域的概率等于该区域面积。令 $A_1,B_1$ 是 $R$ 内的矩形，面积不为 0 或 1。令 $A$ 为随机点落在 $A_1$ 中的事件，$B$ 为随机点落在 $B_1$ 中的事件。给出 $A,B$ 独立的几何描述，并分别给出一个独立例子和一个不独立例子。
    \item 证明若 $A,B$ 独立，则
    \[
      P(A\cup B)=P(A)+P(B)-P(A)P(B)=1-P(A^c)P(B^c).
    \]
  \end{enumerate}

  \item (S) Arby 有一个信念系统，对每个事件 $A$ 赋予一个 $0$ 到 $1$ 之间的数 $P_{\text{Arby}}(A)$。这表示 Arby 对 $A$ 发生可能性的信念程度。对任意事件 $A$，Arby 愿意支付 $1000P_{\text{Arby}}(A)$ 美元购买一种凭证：若 $A$ 发生，该凭证可兑换 1000 美元；若 $A$ 不发生，则无价值。同样，Arby 也愿意按同样价格出售这种凭证。

  假设 Arby 拒绝接受概率公理。特别地，存在两个不交事件 $A,B$，使得
  \[
    P_{\text{Arby}}(A\cup B)\ne P_{\text{Arby}}(A)+P_{\text{Arby}}(B).
  \]
  给出一组 Arby 愿意进行的交易，保证 Arby 亏钱。
\end{enumerate}

\subsection*{容斥原理}

\begin{enumerate}
  \setcounter{enumi}{48}
  \item 掷一枚公平骰子 $n$ 次。至少有一个点数从未出现的概率是多少？

  \item (S) 一名牌手从充分洗匀的标准牌中得到 13 张手牌。求这手牌至少缺少一种花色的概率。

  \item (S) 对 7 个人组成的群体，假设四季（冬、春、夏、秋）作为生日季节等可能，求四季都至少出现一次的概率。

  \item 某班有 20 名学生，周一和周三在一个恰好有 20 个座位的教室上课。某周所有学生两天都出席。两天中，学生都完全随机选座（每座一人）。求没有人两天坐同一座位的概率。

  \item Fred 需要为某网站选择密码。假设他选择 8 个字符的密码，合法字符包括小写字母 a 到 z、大写字母 A 到 Z、以及数字 0 到 9。
  \begin{enumerate}
    \item 若要求密码中至少有一个小写字母，有多少种可能？
    \item 若要求至少有一个小写字母且至少有一个大写字母，有多少种可能？
    \item 若要求至少有一个小写字母、至少有一个大写字母且至少有一个数字，有多少种可能？
  \end{enumerate}

  \item (S) Alice 就读于一所小学院，每门课每周只上一次。她在 30 门互不冲突的课中选择，其中周一到周五每天各有 6 门可选。相信随机性的善意，她决定从 30 门课中随机选择 7 门注册，所有选择等可能。求她周一到周五每天都有课的概率。

  \item 一个俱乐部有 10 名高年级学生、12 名低年级学生和 15 名二年级学生。随机选出一个大小为 5 的组织委员会，所有大小为 5 的子集等可能。
  \begin{enumerate}
    \item 求委员会中恰好有 3 名二年级学生的概率。
    \item 求委员会中高年级、低年级、二年级至少各有一名代表的概率。
  \end{enumerate}
\end{enumerate}

\subsection*{综合练习}

\begin{enumerate}
  \setcounter{enumi}{55}
  \item 对每一小题，判断空格应填 $=,<,>$ 中哪一个，并给出清楚解释。在 (a) 和 (b) 中，顺序不重要。
  \begin{enumerate}
    \item 从 10 人中选 5 人的方式数 \underline{\hspace{1cm}} 从 10 人中选 6 人的方式数。
    \item 把 10 人分成两个 5 人队的方式数 \underline{\hspace{1cm}} 把 10 人分成一个 6 人队和一个 4 人队的方式数。
    \item 3 人组中 3 人都出生在 1 月 1 日的概率 \underline{\hspace{1cm}} 3 人组中 1 人出生在 1 月 1 日、1 人出生在 1 月 2 日、1 人出生在 1 月 3 日的概率。
    \item Martin 和 Gale 玩抛硬币游戏，连续抛一枚公平硬币，直到出现模式 $HH$ 或模式 $TH$。当且仅当 $HH$ 先于 $TH$ 出现时 Martin 获胜。Martin 获胜的概率 \underline{\hspace{1cm}} $1/2$。
  \end{enumerate}

  \item 尝试本题前请深呼吸。假设 Caesar 临终时呼出的最后一口气含有 $10^{22}$ 个分子，大气中共有 $10^{44}$ 个分子。假设如今大气中的分子与 Caesar 生前大气中的分子相同，且在约 2000 年中已经完全随机散布。又假设 Caesar 的最后一口气是无放回抽样，而你刚吸入的一口气是有放回抽样。求你刚吸入的一口气中至少有一个分子也属于 Caesar 最后一口气的概率，并用 $e$ 给出简单近似。

  \item 某检验员检查 12 个小部件，发现恰好 3 个有缺陷。不幸的是，这些部件随后混在一起，检验员必须逐个测试以再次找出 3 个缺陷件。
  \begin{enumerate}
    \item 求检验员现在至少需要测试 9 个部件的概率。
    \item 求检验员现在至少需要测试 10 个部件的概率。
  \end{enumerate}

  \item 有 15 块巧克力和 10 个孩子。在以下情形中，巧克力分给孩子有多少种方式？
  \begin{enumerate}
    \item 巧克力可互换。
    \item 巧克力可互换，且每个孩子至少得到一块。
    \item 巧克力不可互换（哪一块给谁是重要的）。
    \item 巧克力不可互换，且每个孩子至少得到一块。提示：考虑把巧克力随机分给孩子并用容斥原理。
  \end{enumerate}

  \item 给定 $n\ge2$ 个互不相同的数 $(a_1,a_2,\ldots,a_n)$，bootstrap 样本是从这些 $a_j$ 中等概率有放回抽样形成的序列 $(x_1,x_2,\ldots,x_n)$。
  \begin{enumerate}
    \item 对 $(a_1,\ldots,a_n)$，可能的 bootstrap 样本有多少个？
    \item 若顺序不重要，只关心每个 $a_j$ 被选中多少次，可能的 bootstrap 样本有多少个？
    \item 随机选择一个 bootstrap 样本。证明 (b) 中的无序 bootstrap 样本并非等可能。找出一个概率尽可能大的无序样本 $\mathbf b_1$ 和一个概率尽可能小的无序样本 $\mathbf b_2$。令 $p_1$、$p_2$ 分别为得到这两个特定无序样本的概率，求 $p_1/p_2$。
  \end{enumerate}

  \item (S) 100 名乘客排队登上有 100 个座位的飞机，每个座位指定给一名乘客。队伍中第一名乘客随机坐到一个座位上。之后每名乘客若自己的座位可用就坐自己的座位，否则随机坐到一个可用座位上。求最后一名乘客坐到自己指定座位的概率。提示：把第 $j$ 名乘客的座位称为“座位 $j$”。当最后一名乘客上飞机时，哪些座位可能还空着？这些可能性的概率各是多少？

  \item 在生日问题中，我们假设一年 365 天等可能（且排除 2 月 29 日）。现实中有些日子作为生日略微更可能。令 $\mathbf p=(p_1,p_2,\ldots,p_{365})$ 为生日概率向量，其中 $p_j$ 是生日落在一年中第 $j$ 天的概率。第 $k$ 个初等对称多项式定义为
  \[
    e_k(x_1,\ldots,x_n)
    =\sum_{1\le j_1<j_2<\cdots<j_k\le n}x_{j_1}\cdots x_{j_k}.
  \]
  令 $k\ge2$ 为人数。
  \begin{enumerate}
    \item 用 $\mathbf p$ 和初等对称多项式表示至少有一次生日匹配的概率。
    \item 通过考虑简单和极端情形，直观解释为什么当所有 $j$ 都满足 $p_j=1/365$ 时，至少有一次生日匹配的概率最小。
    \item 算术--几何平均不等式说明对 $x,y\ge0$，
    \[
      \frac{x+y}{2}\ge\sqrt{xy}.
    \]
    定义 $\mathbf r=(r_1,\ldots,r_{365})$，其中 $r_1=r_2=(p_1+p_2)/2$，且 $r_j=p_j$ 对 $3\le j\le365$ 成立。利用算术--几何平均不等式和恒等式
    \[
      e_k(x_1,\ldots,x_n)
      =x_1x_2e_{k-2}(x_3,\ldots,x_n)
      +(x_1+x_2)e_{k-1}(x_3,\ldots,x_n)
      +e_k(x_3,\ldots,x_n)
    \]
    证明
    \[
      P(\text{至少一次生日匹配}\mid\mathbf p)
      \ge
      P(\text{至少一次生日匹配}\mid\mathbf r),
    \]
    且当 $\mathbf p\ne\mathbf r$ 时严格不等。用它说明使至少一次生日匹配概率最小的 $\mathbf p$ 满足所有 $p_j=1/365$。
  \end{enumerate}
\end{enumerate}

\end{document}
